% Shared preamble for report.tex and slides.tex.
% Compile both with XeLaTeX: latexmk -xelatex <file>.tex
%
% Visual identity, chosen to be distinct from the Discrete Mathematics template:
%   text     TeX Gyre Pagella (a Palatino), wider and rounder than Times
%   headings TeX Gyre Heros (a Helvetica), sans, to separate structure from prose
%   maths    TeX Gyre Pagella Math through unicode-math, so the digits inside
%            formulas match the digits in the surrounding sentence
% All three ship with TeX Live, so nothing extra has to be installed.

\usepackage{fontspec}
\usepackage{polyglossia}
\setmainlanguage{vietnamese}

\setmainfont{texgyrepagella}[
  Extension      = .otf,
  UprightFont    = *-regular,
  BoldFont       = *-bold,
  ItalicFont     = *-italic,
  BoldItalicFont = *-bolditalic,
]
\setsansfont{texgyreheros}[
  Extension      = .otf,
  UprightFont    = *-regular,
  BoldFont       = *-bold,
  ItalicFont     = *-italic,
  BoldItalicFont = *-bolditalic,
  Scale          = 0.92,
]
\setmonofont{texgyrecursor}[
  Extension      = .otf,
  UprightFont    = *-regular,
  BoldFont       = *-bold,
  Scale          = 0.88,
]

\usepackage{amsmath}
\usepackage{amssymb}
\usepackage{amsthm}
\usepackage{mathtools}
\usepackage{unicode-math}
\setmathfont{texgyrepagella-math.otf}

\usepackage{graphicx}
\usepackage{booktabs}
\usepackage{siunitx}
\usepackage{xcolor}

% Vietnamese decimal marker, so \num agrees with the decimals written by hand in
% the prose (1{,}39 rather than 1.39).
\sisetup{output-decimal-marker={,}}

% Slide copies come first on purpose: \resultgraphic must prefer them.
\graphicspath{{figures/}{../results/figures/slides/}{../results/figures/}}

% ---------------------------------------------------------------------------
% Palette. The three accents are the figure colours from src/figures.py, so a
% method named in the text can be tinted to match its curve.
% ---------------------------------------------------------------------------
\definecolor{ink}{HTML}{1A1A1A}
\definecolor{accent}{HTML}{0B5563}     % deep teal, the structural colour
\definecolor{rule}{HTML}{9AA5A8}
\definecolor{cGD}{HTML}{1F77B4}
\definecolor{cSGD}{HTML}{FF7F0E}
\definecolor{cAGD}{HTML}{2CA02C}
\definecolor{cNewton}{HTML}{D62728}

% ---------------------------------------------------------------------------
% Notation, kept identical to KE_HOACH_TRIEN_KHAI.md
% ---------------------------------------------------------------------------
\newcommand{\R}{\mathbb{R}}
\newcommand{\E}{\mathbb{E}}
\newcommand{\wstar}{w^{*}}
\newcommand{\fstar}{f^{*}}
\newcommand{\bigO}{\mathcal{O}}
\newcommand{\Batch}{\mathcal{B}}
\newcommand{\trans}{^{\top}}
\DeclareMathOperator*{\argmin}{arg\,min}
\DeclareMathOperator{\sign}{sign}
\DeclarePairedDelimiter{\norm}{\lVert}{\rVert}
\DeclarePairedDelimiter{\abs}{\lvert}{\rvert}

% ---------------------------------------------------------------------------
% Figures. A stem that has not been generated yet leaves a visible placeholder
% rather than a compilation error, so the report can be drafted while the
% experiments are still running.
% ---------------------------------------------------------------------------
\newcommand{\missingfigbox}[1]{%
  \fcolorbox{rule}{white}{\parbox[c][3.2cm][c]{0.86\linewidth}{\centering
    \ttfamily\small\detokenize{#1}.pdf\par\vspace{4pt}
    \rmfamily\small chưa được sinh ra từ notebook\par}}%
}

% The path is spelled out rather than left to \graphicspath, which lists the
% slide directory first and would otherwise hand the report the wide, reduced
% copies meant for Beamer.
%
% Placement is [tbp] with no 'h': a figure declared between two paragraphs of one
% argument would otherwise be typeset right there and cut the argument in half.
\newcommand{\resultfig}[3]{%
  \begin{figure}[tbp]
    \centering
    \IfFileExists{../results/figures/#1.pdf}
      {\includegraphics[width=0.82\linewidth]{../results/figures/#1.pdf}}
      {\missingfigbox{#1}}
    \caption{#2}
    \label{#3}
  \end{figure}%
}

% The mandatory pair, stacked rather than side by side. The source figures are
% 6 by 4 inches with 10 pt labels; at half the text width that becomes about 5 pt
% and the legends stop being readable, which defeats the purpose of printing two
% panels. Stacking costs most of a page per comparison and keeps both legible.
\newcommand{\panellabel}[1]{{\sffamily\small\color{accent}#1}}
\newcommand{\resultpair}[3]{%
  \begin{figure}[tbp]
    \centering
    \panellabel{(a) theo số vòng lặp}\par\vspace{2pt}
    \IfFileExists{../results/figures/#1_iters.pdf}
      {\includegraphics[width=0.8\linewidth]{../results/figures/#1_iters.pdf}}
      {\missingfigbox{#1_iters}}
    \par\vspace{10pt}
    \panellabel{(b) theo thời gian chạy}\par\vspace{2pt}
    \IfFileExists{../results/figures/#1_time.pdf}
      {\includegraphics[width=0.8\linewidth]{../results/figures/#1_time.pdf}}
      {\missingfigbox{#1_time}}
    \caption{#2}
    \label{#3}
  \end{figure}%
}

% Slides: figures are 6 by 4 inches, so a full-width copy is too tall for a frame
% that also carries text. Height is capped as a fraction of the text height and
% the aspect ratio kept; crowded frames pass a smaller fraction.
\newcommand{\resultgraphic}[2][0.72]{%
  \IfFileExists{../results/figures/slides/#2.pdf}
    {\includegraphics[width=\linewidth,height=#1\textheight,keepaspectratio]{../results/figures/slides/#2.pdf}}
    {\IfFileExists{../results/figures/#2.pdf}
      {\includegraphics[width=\linewidth,height=#1\textheight,keepaspectratio]{../results/figures/#2.pdf}}
      {\missingfigbox{#2}}}%
}
